\documentclass[instructions.tex]{subfiles}
\begin{document}
 
\chapter{De la libéralité.}
 
Un chrétien doit être libéral, l'être envers tous, l'être avec prudence.
 
\primo On voit des personnes toujours disposées à recevoir, et qui ne donnent jamais rien : c'est la marque d'une âme basse. Il est plus heureux, dit Jésus-Christ, de donner que de recevoir\footnote{Beatius est magis dare quàm acc pere. Act. 20.} . Ayez soin de votre réputation, elle est plus estimable que les richesses. Or, rien ne soutient plus la réputation que le désintéressement et la libéralité.
 
Qui penser d'un homme riche qui ne donne rien, ou qui ne donne que par des vues intéressées? C'est une âme sordide et mercenaire, dont les sentimens n'ont rien de noble et de grand. Quel plaisir peut avoir dans la société un homme dur et tenant, qui n'est estimé de personne?
 
On voit des gens qui paient exactement, qui donnent ce qu'ils promettent, qui ne font tort à personne, qui se croient exempts de blâme, pourvu qu'ils ne fassent point d'injustice; mais si l'on examine de près leur conduite, on n'y voit ni désintéressement ni générosité. Ils craignent toujours de trop donner, de trop promettre, de se dessaisir. Ils disputent, ils chicanent, ils épluchent, pour avoir tout au plus bas prix. Ils ne se fient à personne, pas même à leurs femmes; ils ont l'attention que l'artisan, l'ouvrier, le marchand, ne gagnent trop avec eux ou ne les trompent. Ils ne font point d'injustice, mais aussi ils ne font point de grace. Défians, difficiles, tracassiers dans leurs familles et dans les affaires, ils ne font que des mécontens. Une telle conduite montre ordinairement plus de lésine et d'avarice que d'économie.
 
\secundo Un chrétien, qui a des sentimens de religion et d'honneur, est libéral envers ses amis, dont il sait cultiver l'amitié par des bienfaits. \textit{Les fréquens et petits presens, dit le Proverbe, entretiennent l'union.} Libéral envers ses ennemis, il les prévient par ses services et sa générosité; libéral envers tout le monde, il se fait un plaisir de prêter, de donner, de s'obliger, de relâcher sans vue d'intérêt. Accommodant, facile dans les affaires, il craint plus de faire des mécontens que de perdre. Libéral, surtout envers ceux qui le servent, il paie largement leurs services, parce qu'il comprend que les sueurs, les peines des artisans, des domestiques et de ceux qui nous servent, valent souvent plus que notre argent, qu'il leur coûte plus de nous servir qu'il ne nous coûte de les payer.
 
Une supérieure dit un jour à la sœur qui achetait les provisions du monastère : Ma sœur, ne chicanez point les pauvres gens en achetant leurs denrées, on ne peut trop payer leurs peines; ils font plusieurs lieues pour apporter des charges dont ils n'ont que peu de chose. Quoique nous soyons pauvres, nous ne mourrons pas de faim; apres tout, il vaudrait mieux que nous eussions faim que ces pauvres gens; ils ont plus de maux que nous.
 
Valère-Maxime dit d'un ancien, que tout ce qu'il avait était à tout le monde; que sa maison était l'asile des misérables; qu'il avait un fonds de bienveillance qui le portait à rendre service à tous. Oh! qu'un chrétien qui a cette qualité est estimable aux yeux des hommes, qu'il est cheri de Dieu\footnote{Hilarem enim datorem diligit Deus. 2.Cor.9.}!
 
\tertio Soyez libéral, mais avec prudence et par raison.
 
Faire des dépenses au-delà de ses forces, dissiper son bien en vanité, en jeux, en festins, en débauches, c'est abus, prodigalité, folie et scandale. Il y a cette différence entre le prodigue et l'avare, que le prodigue dissipe comme s'il ne devait vivre que deux jours, et que l'avare amasse comme s'il ne devait jamais mourir. Le premier manque de prévoyance, et l'autre, de confiance en Dieu. La prodigalité est un excès blâmable, et l'avarice, un défaut encore plus honteux, et la marque d'une âme basse.
 
\section{Résolutions}
 
\primo Instruire que les hommes, nés pour la société, doivent s'entraider et se faire mutuellement du bien, chacun selon ses facultés.
 
\secundo Je ne regarderai pas de trop près mes propres intérêts.
 
\tertio Je m'estimerai heureux quand je pourrai rendre service.
 
\quarto Et je ne craindrai pas d'exercer la générosité envers mes semblables, pourvu que ce ne soit ni au préjudice de mes créanciers, ni contre ce que je dois à mon état, ni pour la ruine de ma famille.
 
\prayer{Mon Dieu, faites-moi la grâce d'être libéral, mais de l'être avec prudence, et de sanctifier les services que je rendrai à mon prochain en les lui rendant dans la vue de vous plaire.}
 
\end{document}
